%% kapa: exam=1 round=30 year=2019 period=2 form=A subject=accounting subject_ko=회계학 seq=16 number=56
%% source: Q-Net 감정평가사 기출문제 (cst00309) artlSeq=5205205 file=extracted/1008/5205205/2019년 제30회 감정평가사 1차시험 문제지/2019년 제30회 감정평가사 1차시험 문제지/2019년도 제30회 감정평가사 1차 2교시 A형.hwp
%% conversion: kordoc 4.24.51 -> tools/gen1.py (auto); review=proofread-sonnet; flags=none
\begin{question}{56}
(주)감평은 20x1년 12월 31일자로 종료되는 회계연도 재무제표의 이사회 승인을 앞두고 있다. 아래의 각 상호 독립된 사건은 재무제표에 반영되어 있지 않지만 보고기간말 이후 발생한 것이다. ‘수정을 요하는 보고기간후사건’을 모두 고른 것은? (단, 주석으로 공시되는 금액은 제외한다.)
\begin{framedtext}
ㄱ. 관계회사의 금융기관 차입에 대해 ￦30,000의 지급보증 약정을 체결하였다. \par
ㄴ. 생산공장에 화재가 발생하여 ￦50,000의 생산설비가 파손되었다. \par
ㄷ. 20x1년 말 현재 피고로 계류중이던 손해배상소송에서 ￦10,000의 손해배상 확정판결을 받았다. \par
ㄹ. 내부규정에 의해 20x1년 말 지급하여야 할 상여금 지급액이 ￦25,000으로 확정되었다. \par
\end{framedtext}
\begin{choices}
  \item ㄱ, ㄴ
  \item ㄱ, ㄷ
  \item ㄴ, ㄹ
  \item ㄷ, ㄹ
  \item ㄴ, ㄷ, ㄹ
\end{choices}
\end{question}
\answer{4}
